% © 2026 Ing. David Jaroš — CC BY-NC-ND 4.0
% Licensed under CC BY-NC-ND 4.0. See LICENSE.md.

\ifdefined\INCLUDEMODE
\else
\documentclass[12pt,a4paper]{article}
\usepackage[a4paper,margin=2.5cm]{geometry}
\usepackage{amsmath,amssymb,amsthm}
\usepackage{mathtools}
\usepackage{booktabs}
\usepackage{xcolor}
\usepackage{hyperref}
\usepackage{mdframed}

\newtheorem{theorem}{Theorem}[section]
\newtheorem{proposition}[theorem]{Proposition}
\newtheorem{lemma}[theorem]{Lemma}
\theoremstyle{definition}
\newtheorem{definition}[theorem]{Definition}
\theoremstyle{remark}
\newtheorem{remark}[theorem]{Remark}

\newmdenv[backgroundcolor=green!15,linecolor=green!60!black,linewidth=1.5pt]{resultbox}
\newmdenv[backgroundcolor=yellow!20,linecolor=orange,linewidth=1.5pt]{gapbox}

\title{GAP-U1 Classification: Invariant Symmetric Bilinear Forms on $\mathbb C\otimes\mathbb H\cong M_2(\mathbb C)$}
\author{Ing.\ David Jaroš}
\date{2026-07-14}

\begin{document}
\maketitle
\fi

\section{Problem Setup}

\begin{definition}[Vector space and real structure]
Let
\[
V := M_2(\mathbb C),
\]
viewed as a real vector space of dimension $8$.
\end{definition}

\begin{definition}[Representation]
Fix the biunitary action
\[
\rho(U,V):A\mapsto UAV,\qquad (U,V)\in U(2)_L\times K,
\]
on $V$, where $K\subseteq U(2)$ is the right symmetry subgroup.
\end{definition}

\begin{definition}[Invariant symmetric bilinear form]
A map $\beta:V\times V\to\mathbb R$ is called invariant symmetric bilinear if:
\begin{enumerate}
\item $\beta$ is real-bilinear,
\item $\beta(A,B)=\beta(B,A)$,
\item $\beta(\rho(U,V)A,\rho(U,V)B)=\beta(A,B)$ for all $(U,V)\in U(2)_L\times K$.
\end{enumerate}
\end{definition}

\section{Classification}

For Hermitian $X\in M_2(\mathbb C)$ define
\[
\beta_X(A,B):=\operatorname{Re}\operatorname{Tr}(A X B^\dagger).
\]

\begin{theorem}[Complete classification for fixed right subgroup $K$]
\label{thm:classification-main}
The invariant symmetric bilinear forms of the above type are exactly
\[
\{\beta_X\mid X=X^\dagger,\;VXV^\dagger=X\ \forall V\in K\}.
\]
Equivalently, they are parametrized by the Hermitian part of the commutant
\[
\operatorname{Comm}(K):=\{X\in M_2(\mathbb C)\mid VXV^\dagger=X\ \forall V\in K\}.
\]
\textbf{Verdict: PROVED.}
\end{theorem}

\begin{proposition}[Full biunitary symmetry $K=U(2)$]
If $K=U(2)$ then $\operatorname{Comm}(K)=\mathbb C\,I$, hence
\[
\beta(A,B)=\lambda\,\operatorname{Re}\operatorname{Tr}(AB^\dagger),\qquad \lambda\in\mathbb R.
\]
So the invariant symmetric bilinear form is unique up to scale.
\textbf{Verdict: PROVED.}
\end{proposition}

\begin{proposition}[Lorentzian operator branch $\eta=\mathrm{diag}(1,-1)$]
Define
\[
K_\eta:=\{V\in U(2)\mid V\eta V^\dagger=\eta\}=U(1)\times U(1).
\]
Then
\[
\operatorname{Comm}(K_\eta)_{\mathrm{Herm}}=\{\alpha I+\beta\eta\mid \alpha,\beta\in\mathbb R\},
\]
so the invariant symmetric family is
\[
\beta_{\alpha,\beta}(A,B)=\operatorname{Re}\operatorname{Tr}\!\big(A(\alpha I+\beta\eta)B^\dagger\big).
\]
Therefore $\operatorname{Re}\operatorname{Tr}(A\eta B^\dagger)$ is one member of a two-parameter invariant family, not unique.
\textbf{Verdict: PROVED.}
\end{proposition}

\begin{resultbox}
\textbf{GAP-U1 classification outcome.}
Under full right $U(2)$ symmetry: unique up to scale.
Under the reduced $K_\eta=U(1)\times U(1)$ symmetry induced by fixed $\eta=\mathrm{diag}(1,-1)$: a complete two-parameter family exists.
Hence uniqueness claims must specify the symmetry group.
\end{resultbox}

\ifdefined\INCLUDEMODE
\else
\end{document}
\fi
